\subsection{Harness generation for LLMs}\label{related:generation}
A growing line of work generates task-specific scaffolding around a frozen
model. TTHE evolves text-to-SQL harnesses with an LLM proposer/judge loop
\citep{tthe2026}; Self-Harness \citep{selfharness2026} shows a model can
improve its own harness; JIT-Agent \citep{jitagent2026} trains a model to
generate a harness per task at inference time; Task-CoEvolve
\citep{coevolve2026} co-evolves tasks and scaffolds. A concurrent 2026
cluster makes harness evolution \emph{behavior-aware from the verification
side}: HarnessLens \citep{harnesslens2026} allocates rollout budget using
behavior-level checks during evolution; HarnessFix \citep{harnessfix2026}
repairs harnesses from trace-grounded diagnoses; AHE \citep{ahe2026}
verifies each proposed edit against self-declared behavioral predictions;
gated semantic quality-diversity search \citep{gatedqd2026} evolves
populations over LLM-assigned pathology descriptors; Harness Handbook
\citep{handbook2026} attacks behavior localization and outer-loop search.

Our work is complementary but distinct in unit of analysis and question.
Table~\ref{tab:nearest} contrasts the two nearest systems. HarnessLens uses
behavior as a \emph{signal inside} an evolution loop to allocate compute;
we ask whether the resulting \emph{population} is behaviorally diverse as an
outcome property, and what generation-side factors cause or repair its
collapse. HarnessBank \citep{gatedqd2026} documents search-space collapse and
maintains a semantic diversity library with gated screening; our conformance
gate does not score semantics at all --- it verifies mechanically, via
counterfactual trace probes, that a declared mechanism executes --- and our
design isolates the gate's causal contribution through a factorial contrast
rather than reporting end-to-end system gains.

\begin{table}[t]
\centering
\small
\begin{tabular}{@{}lccc@{}}
\toprule
 & \textbf{This work} & \textbf{HarnessLens} & \textbf{HarnessBank} \\
\midrule
Unit of study & population outcome matrix & evolution loop & diversity library \\
Measures population collapse & yes & no & search collapse only \\
Gate definition & counterfactual traces & attributable-evidence & semantic screening \\
Gate calibration & pos/neg controls & --- & --- \\
Causal isolation of gate & $2{\times}2$ factorial & no & no \\
Confirmatory split & 9 untouched DBs & --- & --- \\
\bottomrule
\end{tabular}
\caption{Nearest concurrent systems. Behavior-aware verification inside
evolution systems is concurrent independent work; our contribution is the
population-level outcome measurement, instrument calibration, and the
controlled generation-side contrast.}
\label{tab:nearest}
\end{table}

\subsection{Harness conditional utility and routing}\label{related:routing}
Several recent papers establish that harness quality is conditional on the
instance: \citet{nouniversal2026} and the non-monotonicity analysis of
\citet{nonmonotone2026} document that no single harness dominates across
instances. HELIX \citep{helix2026} shows a portfolio of sibling harnesses
beats the best fixed one and lists a harness router as future work; STS
\citep{sts2026} routes over six fixed human-designed scaffolds in a closed
set; GRASP \citep{grasp2026} gates skill edits into a library at development
time; \citet{cfrouting2026} learn counterfactual instance-level routing over
\{no-search, search, abstain\}; \citet{agenticrouting2026} route models
within a harness. All of these consume a population whose members
\emph{actually disagree}; \citet{helix2026} comes closest, evaluating a
full candidate outcome matrix and measuring portfolio complementarity beyond
the best fixed candidate. Our contribution is narrower and different: we
measure whether \emph{automatically generated} populations collapse as an
outcome phenomenon, diagnose why at the trace level, and isolate the
generation-side causes with a controlled factorial.
show it fails for naively generated populations, and test whether
generation-side interventions repair it.

\subsection{Behavioral diversity in LLM populations}\label{related:diversity}
Diversity of LLM outputs has mostly been studied at the sample level
(self-consistency, diverse decoding) or the agent-policy level. The
observation that syntactically distinct programs can exhibit identical
behavior echoes coverage gaps documented in program synthesis and
mutation-testing practice; our contribution is to instantiate and quantify
it for AI-generated harness populations, with a mechanism-level taxonomy.
